\paragraph{Statement A.4.7 (Exact smooth model for this particular sharp arc).}
This statement concerns the actual arc, and uses its arithmetic endpoints.
Put
\[
 a_*=17,\qquad B_*=12\sqrt2,\qquad
 R_*=\operatorname{arcosh}(17)=2\operatorname{arcosh}(3),
 \qquad b=\log\frac3{2\sqrt2}.
\]
Thus $a_*^2-B_*^2=1$ and $\coth R_*=\cosh b$.
For $\Re s>0$, define
\[
 W_s(u)=(a_*+B_*\cosh u)^{-s},\qquad
 H_s(\kappa)=\int_{\mathbb R}W_s(u)e^{i\kappa u}\,du,
 \qquad m_s(\kappa)=H_s(\kappa)/K_{s,\kappa}.
\]
Choose a real eigenbasis also diagonalizing the reflection involution.
Every mode that contributes to $p'_j$ is reflection-even. For those modes,
\begin{align}\label{A:U-smooth-boundary}
 \beta_j([0,b])
 &=\frac12\int_{\mathbb R}W_{s_j}(u)\,\beta_j(du),\\
 Q_{j,I}
 &=\frac{K_{s_j,0}}{2L_0}
       \sum_{\ell\in\mathbb Z}p_{j,\ell}\,
                    m_{s_j}(\kappa_\ell).
 \label{A:U-smooth-Q}
\end{align}
The first integral is the distributional pairing with a smooth,
exponentially decreasing function; the second series converges absolutely.
In particular no truncation at the sharp endpoints occurs in
\eqref{A:U-smooth-Q}.

The multiplier has the exact formula
\begin{align}\label{A:U-Legendre-multiplier}
 H_s(\kappa)
 &=\sqrt{\frac{2\pi}{B_*}}\,
   \frac{\Gamma(s+i\kappa)\Gamma(s-i\kappa)}{\Gamma(s)}
   P_{-1/2+i\kappa}^{\,1/2-s}(a_*/B_*),\\
 m_s(\kappa)
 &=\frac{2^{s-1/2}}{\sqrt{\pi B_*}}\,
   \Gamma\!\left(\frac{s+1+i\kappa}{2}\right)
   \Gamma\!\left(\frac{s+1-i\kappa}{2}\right)
   P_{-1/2+i\kappa}^{\,1/2-s}(a_*/B_*).
 \nonumber
\end{align}
For $s=1/2+it$, $t\geq0$, and $|\kappa|\geq t$,
\begin{equation}\label{A:U-forbidden-exponential}
 |m_s(\kappa)|\ll
 (1+t+|\kappa|)^{1/2}
       e^{-\pi(|\kappa|-t)/4}.
\end{equation}

For each fixed $\delta>0$ there is also a uniform stationary-phase
description on $|\rho|=|\kappa/t|\leq1-\delta$ as $t\to\infty$.
Set
\[
 q(\rho)=\sqrt{B_*^2+\rho^2},\qquad
 \xi(\rho)=
 \log\frac{q(\rho)+a_*\rho}{B_*\sqrt{1-\rho^2}},
\]
and
\[
 S(\rho)=\log(a_*+q(\rho))
       -\tfrac12\log(1-\rho^2)-\rho\,\xi(\rho).
\]
Then
\begin{equation}\label{A:U-stationary-multiplier}
 m_{1/2+it}(\kappa)
 =q(\rho)^{-1/2}(1-\rho^2)^{1/4}
       e^{-itS(\rho)}\{1+O_\delta(t^{-1})\}.
\end{equation}
Its canonical time and tangential displacement obey the exact identities
\begin{equation}\label{A:U-canonical-time}
 S'(\rho)=-\xi(\rho),\qquad
 S(\rho)-\rho S'(\rho)
   =\operatorname{arsinh}\!\bigl(B_*\cosh\xi(\rho)\bigr)
   \geq R_*.
\end{equation}
These formulas and the exponential forbidden-region estimate do not by
themselves assert nonvanishing of an eigenspace contraction.

\begin{proof}
All reflections represent the same nontrivial coset of
$W/\Gamma^+$. Thus an odd eigenfunction vanishes pointwise on every
reflection mirror, in particular on both target geodesics.
Its two target periods vanish. We may therefore restrict to
$W$-even eigenfunctions.

In the upper half-plane model with source axis $i e^u$, the reflections
in $z=0$ and $t=0$ act on its real boundary as $x\mapsto-x$ and
$x\mapsto1/x$, respectively. Consequently the boundary distribution
$T_j$ in Statement A.4.6 satisfies $T_j(-dx)=T_j(dx)$, and
$\beta_j$ is even in $u$. The element
\[
 U=S_4S_1,\qquad U(x)=\frac{3x+2\sqrt2}{2\sqrt2 x+3}
\]
belongs to $\Gamma^+$, fixes $x=1$, and maps $\infty$ to $e^b$.
The boundary transformation law, obtained directly from the Poisson
kernel and uniqueness of its distributional boundary value, is
$T_j(U\,dx)=|U'(x)|^{s_j}T_j(dx)$.
On the positive boundary, the two equal signs of $x$ give
$\beta_j(du)=2(2e^u)^{-s_j}T_j(d(e^u))$.
Changing variables $x=U(y)$ on the interval $[1,e^b]$ therefore gives
\[
 \beta_j([0,b])
 =\int_0^\infty
 \left(\frac{e^u}
 {(3e^u+2\sqrt2)(2\sqrt2e^u+3)}\right)^{s_j}
 \beta_j(du)
 =\int_0^\infty W_{s_j}(u)\,\beta_j(du).
\]
Evenness gives \eqref{A:U-smooth-boundary}. The change of variables
is first made away from the endpoints. The boundary primitive is
$1/2$-H\"older, so its interval masses are well-defined; the tail
of the transformed pairing tends to zero because $W_s$ and all its
derivatives decrease exponentially for fixed $\Re s>0$.
This justifies passage to the endpoints.
Periodize $W_s$ with period $L_0$ and use
$p_{j,\ell}=K_{s_j,\kappa_\ell}\beta_{j,\ell}$.
The smooth periodic kernel has rapidly decreasing Fourier coefficients,
which proves \eqref{A:U-smooth-Q}; the stronger convergence assertion
also follows from \eqref{A:U-forbidden-exponential}.

For the special-function identity, take
$x=a_*/B_*>1$, $\nu=-1/2+i\kappa$, and $\mu=s-1/2$
in DLMF \cite[equation (14.12.4)]{DLMF}. Its conditions
$\Re(\mu-\nu)>0$ and $\nu+\mu\notin\{-1,-2,\ldots\}$
hold for $\Re s>0$ and real $\kappa$.
Since $x^2-1=B_*^{-2}$, its integral is exactly one half
of $B_*^sH_s(\kappa)$. This gives the first formula.
The gamma duplication formula gives the second.
For $\kappa=0$, equation (14.12.6), with the conversion
from Olver's $\mathbf Q$ to the standard $Q$, gives
$H_s(0)=2Q_{s-1}(17)$. For example
$H_1(0)=2Q_0(17)=2b$, verifying the constant-mode normalization.

For the uniform bound, the analytic continuation of $W_s(u)$
has no singularity in $|\Im u|<\pi$: its first zeros are
at $\pm b\pm i\pi$.
Shift the contour to $\Im u=3\pi/4$ when $\kappa\geq0$,
and to $\Im u=-3\pi/4$ when $\kappa<0$.
For $v=3\pi/4$ and real $u$,
\[
 |a_*+B_*\cosh(u+iv)|^2
   =288\cosh^2u-408\cosh u+145
   \geq25\cosh^2u.
\]
Also $|\arg(a_*+B_*\cosh(u+iv))|\leq v$.
For $u\geq0$ the imaginary part is nonnegative, and
the imaginary part after multiplication by $e^{-iv}$ equals
$-\sin v\,(a_*+B_*e^{-u}\cos v)<0$;
the other signs follow by conjugation. Hence
\[
 |H_{1/2+it}(\kappa)|
 \leq 5^{-1/2}K_{1/2,0}
       e^{-3\pi(|\kappa|-t)/4}.
\]
The vertical connecting segments vanish by exponential decay in
$|\Re u|$. Uniform two-sided Stirling estimates give
\[
 |K_{1/2+it,\kappa}|
 \asymp
 [(1+|t+\kappa|)(1+|t-\kappa|)]^{-1/4}
 e^{-\frac\pi4(|t+\kappa|+|t-\kappa|-2t)}.
\]
Combining the two estimates proves
\eqref{A:U-forbidden-exponential}.

Finally apply stationary phase separately to $H$ and $K$.
The phases are $\rho u-\log(a_*+B_*\cosh u)$ and
$\rho u-\log\cosh u$.
Their unique stationary points are
\[
 u_W=\log\frac{q(\rho)+a_*\rho}{B_*(1-\rho)},
 \qquad u_K=\operatorname{artanh}\rho .
\]
On $|\rho|\leq1-\delta$ they stay in a fixed compact set and
are uniformly nondegenerate. Outside a larger fixed compact set
the phase derivatives are uniformly separated from zero;
integration by parts treats these exponentially decreasing tails.
At the stationary points, the amplitude divided by the square root
of the absolute second phase derivative equals $q(\rho)^{-1/2}$
for $H$, and $(1-\rho^2)^{-1/4}$ for $K$.
The two stationary-phase factors have the same sign and cancel
in the quotient. This proves \eqref{A:U-stationary-multiplier}.
The formulas for $S'$ and the canonical time follow by differentiation.
Explicitly
$\cosh\xi=q/(B_*\sqrt{1-\rho^2})$ and
$\sinh\xi=a_*\rho/(B_*\sqrt{1-\rho^2})$; substitution gives
\eqref{A:U-canonical-time}.
\end{proof}

\paragraph{Scope of this additional identity.}
The axis of $U$ is the mirror $t=0$, whose normal has norm $-3$.
It is not either norm-$-2$ target.
The two roots of the smooth kernel, $-e^{\pm b}$, are endpoints
of $U^{-1}g_0$, a translate of the source geodesic.
The positive boundary component used in the integral is not a
full period fundamental interval for that axis. Thus the identity
does not turn $Q_{j,I}$ into a positive scalar multiple of $p'_j$.
The nonzero complete-eigenspace coefficient required in
Statement A.5.6 remains a separate question.
